% cap_06_discussao.tex  –  Capítulo 6
% Discussão: SHPN como Formalismo Unificador
\chapter{Discussão: SHPN como Formalismo Unificador}
\label{ch:discussao}

Este capítulo discute as Redes de Petri Hierárquicas de Sinais como extensão
da teoria de Redes de Petri para modelagem de sistemas biológicos onde vias
metabólicas e redes regulatórias devem ser expressas sob um único formalismo.
A contribuição teórica evoluiu para uma \textbf{plataforma computacional de
experimentação \textit{in silico}}, habilitando investigação quantitativa em
biologia de sistemas, farmacocinética, biologia sintética e descoberta de
fármacos.
A discussão examina como essa unificação teórica habilita a predição
quantitativa de limiares de decisão, contextualiza o formalismo dentro
da teoria existente de Redes de Petri, explora quais questões de pesquisa a
plataforma computacional torna tratáveis e identifica problemas teóricos
abertos.

\section{Contribuição Teórica Central}
\label{sec:contribuicao-central}

As Redes de Petri Hierárquicas de Sinais estendem a teoria clássica de Redes
de Petri introduzindo três construtos formais que habilitam modelagem unificada
de vias metabólicas e redes regulatórias: a designação de lugares de sinal
($\Psi \subset P$), os arcos de fluxo de sinal
($F_s \subseteq (\Psi \times T) \cup (T \times \Psi)$)
e as funções de camada hierárquica ($\lambda: \Psi \to \mathbb{N}$).
Esses construtos resolvem uma limitação fundamental nos formalismos existentes
— a incapacidade de modelar controle regulatório por consumo molecular
simultaneamente à transferência de massa metabólica.

\subsection{Unificando Metabolismo e Regulação}

Redes de Petri biológicas clássicas representam vias metabólicas por lugares
(espécies moleculares) e transições (reações) conectados por arcos normais com
pesos estequiométricos.
Arcos de teste estendem esse modelo para catálise por leituras não consumidoras.
Contudo, redes regulatórias operam com semântica diferente: fatores de
transcrição habilitam a expressão gênica enquanto são consumidos em complexos
regulatórios; o ATP fornece energia para biossíntese enquanto a acumulação de Spo0A${\sim}$P
controla programas de desenvolvimento~\cite{Fujita2005}; e moléculas de
sinalização propagam informação por limiares de concentração enquanto são
degradadas ou sequestradas.

O desafio reside em que a mesma molécula participa de ambas as semânticas
simultaneamente — a dualidade metabólito-sinal motivada no
Capítulo~\ref{ch:introducao}.
Abordagens anteriores tratavam esses papéis separadamente, exigindo dois
paradigmas de modelagem distintos.
As SHPN unificam ambos por meio da partição de lugares $\Psi \subset P$: um
mesmo lugar $p \in \Psi$ pode ter arcos normais de entrada representando
produção metabólica e arcos de fluxo de sinal representando consumo regulatório,
formalizando matematicamente a participação dual.

Essa unificação torna três questões de pesquisa anteriormente intratáveis
formalmente expressas.
Em primeiro lugar, a \textbf{predição de limiar a partir de primeiros
princípios}: quando um estado de transição metabólica desencadeia uma
decisão regulatória irreversível?
Formalismos clássicos tratam limiares como parâmetros ajustados; a hierarquia
de sinais os deriva por análise de alcançabilidade, transformando a
identificação de limiar de ajuste de parâmetro em computação formal.
Em segundo lugar, a \textbf{análise composicional de acoplamento
metabólico-regulatório}: vias que compartilham intermediários metabólicos mas
possuem sinais regulatórios distintos ($\Psi_i \cap \Psi_j = \emptyset$) exibem
semântica composicional que habilita modelagem de escala genômica por
decomposição em módulos regulatórios independentes.
Em terceiro lugar, a \textbf{verificação de arquitetura de controle}: a preempção
de camada (Teorema~\ref{thm:preempcao}) impõe controle assimétrico —
sinais de camadas superiores controlam transições de camadas inferiores, mas
camadas inferiores não podem habilitar camadas superiores — tornando a
precedência biológica formalmente verificável.

\section{Posicionamento no Panorama da Teoria de Redes de Petri}
\label{sec:posicionamento-teoria}

A análise de lacunas do Capítulo~\ref{ch:fundamentos} identificou seis
limitações das extensões clássicas de Redes de Petri para modelagem biológica
multi-escala.
As SHPN estendem as Redes Híbridas com três construtos que endereçam essas
lacunas, preservando compatibilidade retroativa ($\Psi = \emptyset$ recupera a
semântica clássica).

A distinção teórica central é a \textbf{estratificação de semântica de arco}
em três categorias --- arcos normais (estequiometria), arcos de fluxo de sinal
(consumo regulatório unidirecional) e arcos de teste (catálise não consumidora)
--- que resolve uma limitação de expressividade: redes clássicas não distinguem
catálise, participação metabólica e controle regulatório.

A função de camada $\lambda$ introduz \textbf{arquitetura de controle formal}
com fluxo assimétrico e preempção hierárquica
(Teorema~\ref{thm:preempcao}), formalizando o princípio de que processos
essenciais (metabolismo) preemptam processos opcionais (diferenciação).
Nenhuma extensão clássica fornece semântica de precedência de controle
comparável~\cite{Simao2025HierarchicalPreemption}, e a verificação de conformidade reduz-se a travessia de grafo
em tempo polinomial.

A semântica de consumo transforma a predição de limiar: em lugar de codificar
limiares como parâmetros de Hill ajustados, o formalismo deriva
$\theta_{\text{eff}}$ analiticamente a partir de $\Gamma$ e identifica o ponto
de comprometimento sistêmico por análise de alcançabilidade
(Capítulo~\ref{ch:validacao-bacillus}).
A \textbf{teoria de independência fraca}
(Teorema~\ref{thm:independencia-fraca}) habilita decomposição composicional,
reduzindo a complexidade de $\mathcal{O}(\bar{M}^n)$ para
$\mathcal{O}(k \cdot \bar{M}^{n/k})$~\cite{Simao2025WeakIndependence,Simao2026QuorumSensing}.

Em síntese, as SHPN quebram três premissas que permeiam todos os formalismos
clássicos de Redes de Petri:

\begin{center}\small
\begin{tabular}{@{}p{3.8cm}p{3.8cm}p{4.5cm}@{}}
\toprule
Premissa clássica & SHPN & Consequência \\
\midrule
Habilitação é plana (verificação local apenas) &
  Habilitação é estratificada ($G_E \wedge$ Preempção) &
  Controle hierárquico representável sem conflação de massa e informação \\
\addlinespace
Todos os arcos são equivalentes (peso é peso) &
  Arcos $F_s$ carregam $\Gamma=(K,n,\varepsilon)$ &
  Limiar molecular, calibrável e cooperativo --- não um inteiro estático \\
\addlinespace
Alcançabilidade é estática a partir de $M_0$ &
  Arcos $F_s$ consumtivos auto-contraem o conjunto alcançável &
  Irreversibilidade estrutural (bacia consumida) --- um teorema, não um
  argumento energético \\
\bottomrule
\end{tabular}
\end{center}

A terceira quebra --- a \textbf{auto-contração do conjunto alcançável} --- é a
contribuição mais profunda em relação às redes clássicas: em qualquer Rede de
Petri clássica é possível reverter um disparo reinjetando tokens.  Nas SHPN,
quando esses tokens pertenciam ao sinal regulatório consumido por $F_s$,
nenhuma injeção externa recupera o contexto de habilitação, porque
$\theta_\text{eff}$ foi computado a partir do estado vivo $M(\psi)$ que o
próprio disparo reduziu.

\section{Questões de Pesquisa Tornadas Tratáveis}
\label{sec:questoes-trataveis}

As SHPN abrem direções de pesquisa anteriormente intratáveis sob métodos formais
existentes.

\subsection{Predição Quantitativa de Limiar para Decisões Biológicas}

A questão central é: em qual estado metabólico uma célula se compromete
irreversivelmente com um programa de desenvolvimento?
Abordagens clássicas tratam os limiares de decisão como parâmetros
medidos experimentalmente e encaixados retroativamente nos modelos.
A hierarquia de sinais inverte isso: dada a topologia e a estequiometria da
rede, o formalismo \emph{prediz} limiares de decisão por análise de
alcançabilidade que computa onde o consumo de sinal cria fronteiras de bacia.

Essa capacidade habilita três aplicações.
\textbf{Paisagens de limiar de escala genômica:} para redes metabólicas com
múltiplos pontos de decisão regulatória, o formalismo habilita o cômputo de
paisagens de limiar multidimensionais que mapeiam estados metabólicos para
destinos de desenvolvimento comprometidos.
\textbf{Projeto de circuitos de biologia sintética:} engenharia de circuitos
sintéticos requer predição de limiares de chave biestável, estabilidade de
estado de memória e confiabilidade de decisão; a hierarquia de sinais fornece
estrutura formal para projeto orientado por limiar.
\textbf{Triagem de alvo terapêutico:} intervenção farmacológica frequentemente
visa deslocar limiares de decisão (indução de apoptose em células
cancerosas, promoção de diferenciação em células-tronco); a plataforma habilita
triagem \textit{in silico} que prediz como perturbações induzidas por fármacos
alteram as paisagens de limiar.

\subsection{Análise Composicional de Crosstalk Regulatório}

A segunda questão é: quando vias biológicas que compartilham intermediários
metabólicos podem ser analisadas independentemente?
A teoria de independência fraca das SHPN estabelece condições formais para
decomposição de vias.
Vias com $\Psi_{\pi_1} \cap \Psi_{\pi_2} = \emptyset$ podem ser analisadas
separadamente mesmo compartilhando lugares metabólicos.
A interseção não vazia $\Psi_{\pi_1} \cap \Psi_{\pi_2} \neq \emptyset$
identifica moléculas de \textit{crosstalk}, tornando a predição de interferência
regulatória uma computação em teoria de grafos em lugar de observação baseada
em simulação.
Para biologia sintética, a independência fraca fornece critério formal de
modularidade: módulos com dependências de sinal disjuntas compõem-se
previsivelmente, abordando o longo desafio de dependência de contexto em
circuitos projetados.

\subsection{Verificação de Arquitetura de Controle Hierárquico}

A terceira questão é: uma proposta de rede biológica respeita a precedência do
fluxo de controle em que o metabolismo controla a regulação, mas não o inverso?
As SHPN tornam a arquitetura de controle uma propriedade verificável: verificar
se os predicados de habilitação referenciam apenas sinais da camada
apropriada é computacionalmente eficiente e formalmente rigoroso.
A verificação inversa — inferir camadas hierárquicas a partir de dados de
perturbação experimental — também se torna possível: observar quais
perturbações controlam quais processos revela a estrutura de camadas por
correspondência com a topologia formal.

\section{Evolução para Plataforma Computacional}
\label{sec:evolucao-plataforma}

A unificação das SHPN de modelagem metabólica e regulatória sob um único
formalismo evoluiu além da contribuição teórica para uma \textbf{plataforma
computacional de experimentação \textit{in silico}} em domínios diversos.

Na \textbf{biologia de sistemas}, a plataforma fornece estrutura unificada que
integra dados multi-ômicos (metabolômica, transcriptômica, proteômica) por
representação computacional única abrangendo metabolismo e regulação, tornando
o teste quantitativo de hipóteses tratável.
Na \textbf{farmacocinética e descoberta de fármacos}, a ação de fármacos envolve
transformação metabólica acoplada a resposta regulatória; a plataforma modela
ambas simultaneamente, habilitando triagem \textit{in silico} que prediz tanto
eficácia (cruzamento de limiar) quanto toxicidade (efeitos metabólicos
inespecíficos)~\cite{Simao2026EmergentThreshold}.
Na \textbf{biologia sintética}, circuitos acoplam metabolismo (disponibilidade
de recursos, custos energéticos) a lógica regulatória projetada; a predição de
limiar orienta a seleção de moléculas de sinal e a análise de independência
fraca identifica fronteiras de módulos.

\section{Problemas Teóricos em Aberto}
\label{sec:problemas-abertos}

\subsection{Otimização do Algoritmo de Alcançabilidade}

A teoria de independência fraca reduz a complexidade de análise de exponencial a
polinomial decompondo as redes em módulos independentes.
No entanto, a implementação computacional pode ser otimizada adicionalmente.
Um problema em aberto é se a estrutura hierárquica de camadas pode habilitar
computação paralela de alcançabilidade: se as interações entre vias formam um
grafo de dependência acíclico, a análise paralela de alcançabilidade pode
atingir aceleração quase linear.
Outro problema em aberto é se a aproximação limitada por camada pode computar
limites de limiar eficientemente: se os sinais de camadas inferiores dominam a
determinação de limiar, analisar apenas as camadas de baixo pode aproximar os
limiares da rede completa com limites de erro prováveis.

\subsection{Conexão com Transições de Fase para Estados Absorventes}
\label{subsec:estados-absorventes}

Um estado $s$ é \emph{absorvente} no sentido da física estatística se
$P(s \to s) = 1$ --- uma vez atingido, nenhuma transição de saída existe.
Sistemas com estados absorventes exibem \emph{transições de fase para estados
aborventes}~\parencite{Hinrichsen2000}: abaixo de um parâmetro crítico
$\lambda_c$ o sistema sempre atinge o estado absorvente; acima de $\lambda_c$
a atividade é sustentada indefinidamente; exatamente em $\lambda_c$ emergem
distribuições de lei de potência e expoentes críticos universais (classe de
universalidade de percolação dirigida --- Janssen 1981, Grassberger 1982).

O estado comprometido em SHPN possui a estrutura de um estado absorvente: após
o consumo dos tokens de sinal por $F_s$, não existem transições de saída ---
estruturalmente, não energeticamente (Observação~\ref{rem:bacia-consumida}).
O mapeamento entre as duas teorias é:

\begin{center}\small
\begin{tabular}{@{}ll@{}}
\toprule
Teoria de estados absorventes & SHPN \\
\midrule
Estado absorvente & Marcação comprometida $M^*$ pós-cascata \\
Parâmetro crítico $\lambda_c$ & $\theta_\text{eff}$ \\
Regime sub-crítico & $M(\psi) < \theta_\text{eff}$: retorna ao atrator vegetativo \\
Regime super-crítico & $M(\psi) \geq \theta_\text{eff}$: absorvido no estado comprometido \\
Parâmetro de ordem & $\varepsilon$ de $\Gamma$ \\
\bottomrule
\end{tabular}
\end{center}

O componente $\varepsilon$ de $\Gamma$ atua como \emph{parâmetro de ordem da
classe de irreversibilidade}: quando $\varepsilon \to 0$ o estado comprometido
é verdadeiramente absorvente --- supressão completa sem caminho reverso
(esporulação, apoptose).  Quando $\varepsilon > 0$ o comprometimento é
metaestável, com atividade reversa residual proporcional a $\varepsilon$
(diferenciação de células-tronco reprogramável por Yamanaka~\parencite{Takahashi2006}).

Um problema em aberto é determinar se o comprometimento em SHPN pertence à
classe de universalidade de percolação dirigida: trajetórias de $\tau$-leaping
próximas a $\theta_\text{eff}$ deveriam exibir distribuição de lei de potência
para os tempos de comprometimento e susceptibilidade divergente,
independentemente dos detalhes moleculares do sistema em questão.

\subsection{Fronteiras de Expressividade}

As SHPN modelam o controle regulatório por consumo de sinal baseado em limiar,
o que captura bem decisões mas pode expressar
inadequadamente outros fenômenos biológicos.
\textbf{Gradientes espaciais contínuos:} enquanto o formalismo suporta
compartimentalização espacial discreta, gradientes de concentração contínuos
(difusão de morfogênio, ondas de cálcio) exigem extensões de campo espacial;
integrar a semântica de sinal com extensões de Rede de Petri de reação-difusão
permanece em aberto.
\textbf{Decisões probabilísticas:} algumas decisões de desenvolvimento parecem
estocásticas mesmo sob condições iniciais idênticas; desenvolver teoria de
alcançabilidade probabilística para Redes de Petri de sinal habilitaria o
cômputo de probabilidade de decisão como função do estado metabólico.
\textbf{Restrições temporais:} desenvolvimento envolve sequências temporais e
requisitos de duração; integrar a semântica de sinal com extensões de Rede de
Petri temporal permanece inexplorado.

\subsection{Escopo de Validação}

A validação atual focou na esporulação bacteriana — uma decisão unicelular em
desenvolvimento procariótico.
A regulação gênica eucariótica envolve remodelação da cromatina, modificações
de histona e organização nuclear espacial; se a hierarquia de sinais captura
essas camadas regulatórias adicionais requer validação experimental.
Processos de desenvolvimento multicelular que envolvem comunicação célula-célula
(gradientes de morfogênio, inibição lateral Notch/Delta) requerem extensões
espaciais.
Mecanismos regulatórios não metabólicos — limiares de concentração de morfogênio
(Bicóide em \textit{Drosophila}), mecanotransdução (YAP/TAZ), potencial elétrico
— não foram formalizados; se a semântica de consumo de sinal se aplica
universalmente ou especificamente ao controle metabólico requer validação
adicional.

\section{Síntese do Capítulo}
\label{sec:sintese-discussao}

As SHPN posicionam-se como extensão das Redes Híbridas com três construtos
novos --- lugares de sinal, arcos de fluxo de sinal e camadas hierárquicas ---
que habilitam predição de limiar por análise de alcançabilidade
(Seção~\ref{sec:posicionamento-teoria}), análise composicional por independência
fraca e verificação de arquitetura de controle.
As questões de pesquisa tornadas tratáveis
(Seção~\ref{sec:questoes-trataveis}) e os problemas teóricos em aberto
(Seção~\ref{sec:problemas-abertos}) delineiam o programa de investigação
futuro.
